\scanpage{44}
% archive foliation: 43
\diagram{\begin{tikzcd}&X\arrow[d]\\S&T\arrow[l,"f"']\end{tikzcd}}
$f_*(X/T)$ \qquad $T$ fini sur $S$
\diagram{\begin{array}{llll}
\text{Si}&X/T\text{ est affine},&f_*(X/T)\text{ aussi}.&\\
\text{si}&\text{---}\left\{\begin{array}{l}\text{projectif}\unsure{*|+}\\\text{propre}\end{array}\right.&\text{---}&\left.\vphantom{\begin{array}{l}1\\2\end{array}}\right\}\ \uline{\text{pourvu que}}\ f\ \uline{\text{séparable}}\ (*)\\
\text{si}&\text{quasi-projectif},&\text{---}&
\end{array}}
\margin{\unsure{donner un énoncé} ``général'' \unsure{dans le cas}: \unsure{$f$ étale|$f$ fidèle}.}

Lemme $f_*$ est exact à gauche, \deleted{\unsure{vérifie que le produit}} \inserted{en particulier}
\[
f_*(T/T)=S,\quad f_*(X\times_T Y/T)=f_*(X/T)\times_S f_*(Y/T),\quad f_*(X\times_Z Y/T)=f_*(X/T)\times_{f_*(Z/T)}f_*(Y/T).
\]
Le foncteur $f_*$ transforme imm. ouvertes en imm. \deleted{\unsure{fermées}} \inserted{ouvertes} et \unsure{l'on a une hyp. qui assure}, \deleted{[\unsure{pourvu que}} \inserted{\unsure{finies}}
\[
f_*(X/T)=\bigcup_i f_*(U_i/T),\quad\text{où les }U_i\text{ recouvrent }X]
\]
imm. fermées en imm. fermées. \deleted{\unsure{Les}} \deleted{\unsure{autres assertions},}
[N.B. cette assertion doit pouvoir être \unsure{déduite|établie} par la technique de descente].
\inserted{\unsure{que le $=$ visé est même : à vérifier;}}

\margin{(*) Se voit par descente:\par
Si $Y/S$ est tel que $f^*(Y)/S'$ soit propre resp. projectif, alors $Y/S$ l'est [Dans le deuxième cas, on suppose $T/S$ fini et plat sur $S$ \uline{loc. noeth.} \deleted{\unsure{(1)}}\par
(i) \unsure{faire détail dans le cas actuel} pour voir l'assertion, en se ramenant au projectif type, \deleted{\illegible{2}} travaillant sur $\underline{Z}$ en définitive, il faut prouver ceci:\par
Si un faisceau loc. libre inversible $\uline{L}$ sur $X$ est $T$-ample, alors $f_*(\uline{L})$ [qui est un faisceau loc. libre de rang $d$ sur $f_*(X/T)$ est tel que $\det f_*(\uline{L})$ est $S$-ample --\par
Dans le premier cas, on utilise le lemme de Chow, en notant que si $X'\longrightarrow X$ est surjectif, alors $f_*(X')\longrightarrow f_*(X)$ l'est -- \deleted{\illegible{2}}, on peut pour cela supposer $S=(\operatorname{spec} k)$; et \unsure{préciser par descente} ---\par
Ici il faut $S$ \deleted{\unsure{noeth.}} \inserted{loc.} noeth.\par
\uline{Question} Si $T/S$ est fidèlement plat non fini, et $\unsure{X/S|Y/S}$ est tel que $Y\times_S T/T$ soit propre, est-il vrai que $Y/S$ l'est? -- Oui dans le cas $S$ loc. noeth., car on utilise encore le lemme de Chow.}

De $=$, $f_*$ transforme \deleted{\unsure{applications}} morphismes affines en morphismes affines [en effet, soit $\varphi:X\longrightarrow Y$ affine, \deleted{alors} on est ramené \inserted{par recollement} au cas où $Y$ est affine et $X\subset Y\deleted{\unsure{\otimes_T}} Z\deleted{\unsure{\subset X}}$ est schéma fermé \inserted{$Z/T$ affine}, on se ramène au cas $\unsure{X\supset Y\otimes_T Z|X=Y\otimes_T Z}$ puis on tient des imm. fermées] il transforme morphismes projectifs en morphismes projectifs si $T/S$ est séparable, et en tous cas il transforme morphismes quasi-projectifs en morphismes quasi-projectifs.

[Méthode, du moins si $Y$ est \uline{noeth.} \inserted{\unsure{si on $X_T\longrightarrow Y/T$}} \deleted{\unsure{du}} qu. proj. de type fini sur $T$ donc un faisceau coh. sur $Y$ est quotient d'une puissance de $T$; \unsure{Il vaut} mieux, et il est plus précis de pouvoir assurer que l'on
